% =====================================================================
%  c05.tex —— §6 反例（原书 p.53–54；PDF p.12（§6 起）–p.13（§7 标题前））
%  编号公式：本块无
%  定理类环境：例 6.1、命题 6.2、命题 6.3、注意 6.4
%  图表：1 幅（例 6.1 的态射示意图，tikz-cd 重画，无编号）
%  说明：p.12 顶部两行是 §5 的末段，属 c04；§7 标题属 c06，本块均不含。
% =====================================================================

% ===== 原书 p.53 =====
\section{反例}
\label{sec:c05}

[12] 中构造了问题 4.1 的反例，但那里在次数上含有错误（参见下面的注意 6.4）。以下给出订正后的构造。

\begin{example}
  设极小$\ainf$范畴$s\cC$如下。
  \[
  \begin{tikzcd}[column sep=1.6em, row sep=1em]
    X_0 \arrow[r,"\rho_{01}"]
    & X_1 \arrow[r,"\rho_{12}"]
    & X_2 \arrow[r,"\rho_{23}"]
    & X_3 \arrow[r,"\rho_{34}"]
    & X_4 \arrow[r,"\rho_{45}"]
    & X_5 \arrow[r,"\rho_{56}"]
    & X_6 \arrow[r,"\rho_{67}"]
    & X_7.
  \end{tikzcd}
  \]
  $\cC$的对象为$X_0,X_1,\dots,X_6,X_7$。
  \begin{enumerate}[label=(\roman*),leftmargin=2.8em,labelsep=0.4em,
                    itemsep=2pt,topsep=4pt,parsep=0pt]
  \item 对$i=0,\dots,6$，态射空间$s\cC(X_i,X_{i+1})$设为以次数$0$的基
        $\rho_{i(i+1)}$张成的$1$维向量空间。（至此对应上图。）
  \item 态射空间$s\cC(X_i,X_{i+2})$：当$i=2,3$时，设为以次数$1$的基
        $\rho_{i(i+2)}$与次数$0$的基$\sigma_{i(i+2)}$张成的$2$维向量空间；
        其余的$i$则设$s\cC(X_i,X_{i+2})=0$。
  \item 态射空间$s\cC(X_1,X_4)$、$s\cC(X_3,X_6)$分别设为以基$\rho_{14}$、
        $\rho_{36}$张成的次数$1$的$1$维向量空间。其余$(i\neq1,3)$的情形设
        $s\cC(X_i,X_{i+3})=0$。
  \item 态射空间$s\cC(X_0,X_4)$、$s\cC(X_3,X_7)$分别设为以次数$1$的基
        $\rho_{04}$、$\rho_{37}$张成的$1$维向量空间。其余$(i\neq0,3)$的情形设
        $s\cC(X_i,X_{i+4})=0$。
  \item 对$i=0,1,2$全部设$s\cC(X_i,X_{i+5})=0$。
  \item $i=0,1$两种情形均设$s\cC(X_i,X_{i+6})=0$。
  \item $s\cC(X_0,X_7)$设为以次数$2$的基$\rho_{07}$张成的$1$维向量空间。
  \end{enumerate}
  此外，对每个$i=0,1,\dots,6,7$，$s\cC(X_i,X_i)$设为以恒等态射
  $s(\id_{X_i})$张成的次数$-1$的$1$维向量空间；$i>j$时设
  $s\cC(X_i,X_j)=0$。

  固定某个$\epsilon\in K$，如下定义$\cC$的极小$\ainf$结构（即$m_1=0$）。
  除含恒等元$s(\id_{X_i})$的积$m_2$之外，非平凡的$\ainf$积定义如下：
  \[
  \begin{gathered}
    m_2(\rho_{23},\rho_{34})=\rho_{24},\qquad
    m_2(\rho_{34},\rho_{45})=\rho_{35},\\
    m_2(\rho_{12},\sigma_{24})=\rho_{14},\qquad
    m_2(\sigma_{35},\rho_{56})=\rho_{36}\\
    m_3(\rho_{12},\rho_{23},\rho_{34})=\epsilon\cdot\rho_{14},\qquad
    m_3(\rho_{34},\rho_{45},\rho_{56})=-\epsilon\cdot\rho_{36}\\
    m_3(\rho_{01},\rho_{12},\sigma_{24})=-\rho_{04},\qquad
    m_3(\sigma_{35},\rho_{56},\rho_{67})=-\rho_{37}\\
    m_4(\rho_{04},\rho_{45},\rho_{56},\rho_{67})=\rho_{07}
    =m_4(\rho_{01},\rho_{12},\rho_{23},\rho_{37}),\\
    m_5(\rho_{01},\rho_{14},\rho_{45},\rho_{56},\rho_{67})=\rho_{07}
    =m_5(\rho_{01},\rho_{12},\rho_{23},\rho_{36},\rho_{67}).
  \end{gathered}
  \]

  % ===== 原书 p.54 =====
  此时存在如下非平凡的$\ainf$关系式：
  \[
  \begin{gathered}
    m_5(\rho_{01},m_3(\rho_{12},\rho_{23},\rho_{34}),\rho_{45},\rho_{56},\rho_{67})
    +m_5(\rho_{01},\rho_{12},\rho_{23},m_3(\rho_{34},\rho_{45},\rho_{56}),\rho_{67})=0,\\
    m_5(\rho_{01},m_2(\rho_{12},\sigma_{24}),\rho_{45},\rho_{56},\rho_{67})
    +m_4(m_3(\rho_{01},\rho_{12},\sigma_{24}),\rho_{45},\rho_{56},\rho_{67})=0,\\
    m_5(\rho_{01},\rho_{12},\rho_{23},m_2(\sigma_{35},\rho_{56}),\rho_{67})
    +m_4(\rho_{01},\rho_{12},\rho_{23},m_3(\sigma_{35},\rho_{56},\rho_{67}))=0.
  \end{gathered}
  \]
  记这个极小$\ainf$范畴为$s\cC_\epsilon$。
  \label{exa:c05-6.1}
\end{example}

\transnote{
这里出现的次数都以悬垂范畴$s\cC$为准。由$(s\cC)^{i}(X,Y)=\cC^{i+1}(X,Y)$，
$\cC$中次数为$d$的态射在$s\cC$中次数为$d-1$；所以$\cC$中次数$0$的恒等态射
$s(\id_{X_i})$在$s\cC$中次数为$-1$，以下各$\rho$、$\sigma$的次数也全部按$s\cC$计。
用未悬垂的约定去读这些数字，会整体差$1$。
}

\transnote{
上面的三个关系式分别是$n=7$（第一式）与$n=6$（后两式）的$\ainf$关系式中
残留的项：内层积取$m_3$或$m_2$，外层积相应取$m_5$或$m_4$，其余位置上的积全为
零，因此只剩这两项。逐项核对次数可验证这三个式子都是齐次的。
}

由$\epsilon=0$时的$\cC_0$与$\epsilon\neq0$时的$\cC_\epsilon$组成的对
$(\cC_0,\cC_\epsilon)$给出反例。为简单起见，$\epsilon\neq0$的情形取
$\epsilon=1$考虑。

\transnote{
原文在给这个范畴命名时写作$s\cC_\epsilon$，但在写对$(\cC_0,\cC_\epsilon)$时
又省去$s$，记作$\cC_0$、$\cC_\epsilon$。下文两者指同一个范畴，即
$\cC_\epsilon=s\cC_\epsilon$。
}

\begin{proposition}[{[12, Proposition 4.21]}]
  $\ainf$范畴的对$(\cC_0,\cC_1)$满足定理 5.4 的条件 (i), (ii), (iii)。
  \label{prop:c05-6.2}
\end{proposition}

\begin{proposition}[{[12, Proposition 4.22]}]
  这两个极小$\ainf$范畴$\cC_0$、$\cC_1$互不$\ainf$同型。
  \label{prop:c05-6.3}
\end{proposition}

这两个命题的证明省略，但要点如下。命题 6.2 由下述事实得出：依赖于$\epsilon$的
两个$3$元$\ainf$积$m_3$对扭积（定义 5.1）的贡献，只剩下含于空间$B^1$的部分。
至于命题 6.3，必须脚踏实地地说明不存在$\ainf$同型；不过此时态射空间满足$i>j$时
$s\cC(X_i,X_j)=0$这一性质，即$\cC_\epsilon$属于所谓有序$\ainf$范畴（cf.\ [20]），
因此其证明属于可完成的水平。由这两个命题与定理 4.2，可知这个
$(\cC_0,\cC_1)$是问题 4.1 的反例。

\begin{remark}
  [12] 中给出了这个反例，但那里在态射的次数上，(iii) 的部分差了一个，随之 (vii)
  部分的次数也差了一个。应当按上面这样订正 [12, Example 4.20] 中的这处错误。
  这一订正不影响上面的命题 6.2；据 Deroen Maes 告知[18]，命题 6.3 也能不作改动、
  同样地证明。
  \label{rem:c05-6.4}
\end{remark}

\transnote{
订正后的次数可以这样检验：$m_2$的次数为$1$，故$m_2(\rho_{12},\sigma_{24})
=\rho_{14}$要求$\rho_{14}$的次数为$1$，若按 (iii) 少算一，这个积就不齐次。
同理，$m_4$的次数为$1$，$m_4(\rho_{04},\rho_{45},\rho_{56},\rho_{67})=\rho_{07}$
左端输入次数之和为$1$，合计$2$，故$\rho_{07}$的次数必须是$2$。
}
